% Cours — Programmation orientée objet (BTS CIEL 1re année)
% Diaporama Beamer, palette et encadrés alignés sur les autres supports du dépôt.
% Compilation : pdflatex Cours_POO.tex  (deux passes pour la table des matières)
\documentclass[11pt,aspectratio=169]{beamer}

\usepackage[T1]{fontenc}
\usepackage[utf8]{inputenc}
\usepackage{lmodern}
\usepackage[french]{babel}
\usepackage{xcolor}
\usepackage{booktabs,tabularx,array}
\usepackage{amssymb,amsmath}
\usepackage[most]{tcolorbox}
\usepackage{listings}
\usepackage[strings]{underscore}   % permet d'écrire __init__ sans échappement
\usepackage{tikz}
\usetikzlibrary{positioning,arrows.meta,shapes.geometric,fit,backgrounds}

% ---------------------------------------------------------------- palette
% Reprise exacte de la palette des supports article du dépôt
\definecolor{ipred}{HTML}{C0392B}
\definecolor{cdcbar}{HTML}{A67C3D}
\definecolor{cdcbg}{HTML}{F4F1EA}
\definecolor{qbar}{HTML}{4A6B5D}
\definecolor{qbg}{HTML}{EDF1EF}
\definecolor{corr}{HTML}{1F5FA8}
\definecolor{corrbg}{HTML}{E8F0FA}
\definecolor{grisfoot}{HTML}{828282}
\definecolor{codebg}{HTML}{F7F7F4}
\definecolor{codekw}{HTML}{8E44AD}
\definecolor{codestr}{HTML}{1F7A4D}
\definecolor{codecom}{HTML}{8A8A8A}

% ---------------------------------------------------------------- thème
\usetheme{default}
\usecolortheme{default}
\setbeamertemplate{navigation symbols}{}
\setbeamercolor{structure}{fg=qbar}
\setbeamercolor{frametitle}{fg=qbar,bg=}
\setbeamercolor{title}{fg=qbar}
\setbeamercolor{itemize item}{fg=cdcbar}
\setbeamercolor{itemize subitem}{fg=qbar}
\setbeamercolor{section in toc}{fg=black}
\setbeamerfont{frametitle}{series=\bfseries,size=\large}
\setbeamertemplate{itemize item}{\color{cdcbar}\raisebox{0.15ex}{\footnotesize$\blacktriangleright$}}
\setbeamertemplate{itemize subitem}{\color{qbar}\scriptsize$\bullet$}

% Titre de diapo : filet doré sous le titre, comme les encadrés des polycopiés
\setbeamertemplate{frametitle}{%
  \vspace{1.5mm}%
  \hspace{6mm}\begin{minipage}{\dimexpr\paperwidth-14mm\relax}
    \raggedright\usebeamerfont{frametitle}\color{qbar}\insertframetitle\par
    \vspace{1mm}{\color{cdcbar}\rule{\linewidth}{1.1pt}}
  \end{minipage}\par\vspace{0.5mm}}

% Pied de page sobre, identique en esprit au \fancyfoot des polycopiés
\setbeamertemplate{footline}{%
  \hbox{\begin{beamercolorbox}[wd=\paperwidth,ht=2.4ex,dp=1.2ex,leftskip=6mm,rightskip=6mm]{}%
    {\small\itshape\color{grisfoot}Programmation orientée objet}\hfill
    {\small\color{grisfoot}BTS CIEL 1}\hfill
    {\small\color{grisfoot}\insertframenumber}%
  \end{beamercolorbox}}}

\setbeamertemplate{blocks}[default]

% ---------------------------------------------------------------- code Python
\lstdefinestyle{py}{
  language=Python,
  basicstyle=\ttfamily\scriptsize,
  keywordstyle=\color{codekw}\bfseries,
  stringstyle=\color{codestr},
  commentstyle=\color{codecom}\itshape,
  showstringspaces=false,
  backgroundcolor=\color{codebg},
  frame=single,
  rulecolor=\color{black!25},
  framesep=5pt,
  xleftmargin=4pt,
  xrightmargin=4pt,
  columns=fullflexible,
  keepspaces=true,
  breaklines=true,
  extendedchars=true,
  literate={é}{{\'e}}1 {è}{{\`e}}1 {ê}{{\^e}}1 {à}{{\`a}}1 {ç}{{\c c}}1
           {ù}{{\`u}}1 {û}{{\^u}}1 {î}{{\^\i}}1 {ô}{{\^o}}1 {É}{{\'E}}1 {À}{{\`A}}1 {È}{{\`E}}1 {Ç}{{\c C}}1
           {«}{{\guillemotleft}}1 {»}{{\guillemotright}}1,
  morekeywords={self,True,False,None},
}
\lstset{style=py}
\newcommand{\py}[1]{\texttt{#1}}

% ---------------------------------------------------------------- encadrés
\newtcolorbox{defi}[1][]{enhanced,colback=qbg,colframe=qbar,boxrule=0pt,
  leftrule=3pt,arc=0pt,outer arc=0pt,left=6pt,right=6pt,top=4pt,bottom=4pt,
  title={#1},coltitle=qbar,fonttitle=\bfseries,attach title to upper={\ :\ }}
\newtcolorbox{methode}[1][]{enhanced,colback=cdcbg,colframe=cdcbar,boxrule=0pt,
  leftrule=3pt,arc=0pt,outer arc=0pt,left=6pt,right=6pt,top=4pt,bottom=4pt,
  title={#1},coltitle=cdcbar,fonttitle=\bfseries,detach title,
  before upper={\tcbtitle\par\smallskip}}
\newtcolorbox{piege}{enhanced,colback=red!4,colframe=ipred,boxrule=0pt,
  leftrule=3pt,arc=0pt,outer arc=0pt,left=6pt,right=6pt,top=4pt,bottom=4pt}
\newtcolorbox{exam}{enhanced,colback=corrbg,colframe=corr,boxrule=0pt,
  leftrule=3pt,arc=0pt,outer arc=0pt,left=6pt,right=6pt,top=4pt,bottom=4pt}

\setbeamersize{text margin left=6mm,text margin right=6mm}
\setlength{\parskip}{3pt}

% ---------------------------------------------------------------- styles TikZ
\tikzset{
  boite/.style={draw=black!55,fill=white,rounded corners=1pt,inner xsep=5pt,inner ysep=3pt,font=\small},
  classe/.style={boite,fill=qbg,draw=qbar},
  objet/.style={boite,fill=cdcbg,draw=cdcbar},
  fleche/.style={-{Stealth[length=3mm]},draw=corr,thick},
}

\title{Programmation orientée objet}
\subtitle{Classes, objets et les trois piliers}
\author{}
\date{BTS CIEL --- 1\textsuperscript{re} année}

\begin{document}

% ================================================================= titre
{\setbeamertemplate{footline}{}
\begin{frame}[plain]
  \vfill
  \begin{center}
    {\color{cdcbar}\rule{0.5\linewidth}{1.2pt}}\par\medskip
    {\LARGE\bfseries\color{qbar}Programmation orientée objet\par}\medskip
    {\large Classes, objets et les trois piliers\par}\medskip
    {\color{cdcbar}\rule{0.5\linewidth}{1.2pt}}\par\bigskip
    {\color{grisfoot}BTS CIEL --- 1\textsuperscript{re} année \quad|\quad Python}
  \end{center}
  \vfill
\end{frame}}

\begin{frame}{Plan}
  \tableofcontents
\end{frame}

% ================================================================= 1. l'idée
\section{L'idée de la POO}

\begin{frame}{Qu'est-ce que la POO ? L'analogie}
  \begin{defi}[Idée principale]
    programmer en décrivant des \textbf{objets}, à l'image de ceux de la vie réelle,
    plutôt qu'une simple suite d'instructions.
  \end{defi}
  \medskip
  \textbf{L'analogie de la voiture}
  \begin{itemize}
    \item on écrit une \emph{recette} (un modèle) qui décrit ce qu'est une « voiture » ;
    \item à partir de cette recette, on fabrique autant de voitures concrètes qu'on veut.
  \end{itemize}
  \medskip
  \begin{center}
  \begin{tikzpicture}[node distance=3.2cm]
    \node[classe] (c) {\begin{tabular}{c}\textbf{Vehicule}\\\hline \py{nombre\_roues}\end{tabular}};
    \node[objet,right=5.5cm of c] (o) {\begin{tabular}{c}\textbf{moto\_1}\\\hline \py{nombre\_roues = 2}\end{tabular}};
    \node[above=1mm of c,font=\small\color{grisfoot}] {Classe};
    \node[above=1mm of o,font=\small\color{grisfoot}] {Objet};
    \draw[fleche] ([yshift=3mm]c.east) -- node[above,font=\small]{construire} ([yshift=3mm]o.west);
    \draw[fleche] ([yshift=-3mm]c.east) -- node[below,font=\small]{instancier} ([yshift=-3mm]o.west);
  \end{tikzpicture}
  \end{center}
\end{frame}

% ================================================================= 2. classe
\section{Classe et objet}

\begin{frame}[fragile]{La « recette » $\rightarrow$ la classe}
  \begin{defi}[Classe]
    modèle (ou plan, ou moule) qui définit les \textbf{caractéristiques} et les
    \textbf{comportements} des objets créés à partir d'elle.
  \end{defi}
  \begin{columns}[T,onlytextwidth]
    \begin{column}{0.46\textwidth}
      Une classe définit deux choses :
      \begin{itemize}
        \item les \textbf{attributs} --- ce qu'un objet \emph{a}\\
              {\small couleur, marque, vitesse}
        \item les \textbf{méthodes} --- ce qu'un objet \emph{fait}\\
              {\small \py{rouler()}, \py{freiner()}}
      \end{itemize}
    \end{column}
    \begin{column}{0.52\textwidth}
\begin{lstlisting}
class Voiture:
    # Le constructeur : initialise les attributs
    def __init__(self, couleur, marque):
        self.couleur = couleur
        self.marque  = marque
        self.vitesse = 0

    # Une méthode (action)
    def rouler(self):
        self.vitesse += 10
\end{lstlisting}
    \end{column}
  \end{columns}
\end{frame}

\begin{frame}[fragile,shrink=5]{L'objet $\rightarrow$ l'instance}
  \begin{defi}[Objet]
    exemplaire concret créé à partir d'une classe. On dit que c'est une
    \textbf{instance} de cette classe.
  \end{defi}
  Un objet peut représenter une entité du monde physique ou un concept :
  un individu, une voiture, un nombre complexe, une matrice\dots
\begin{lstlisting}
# Instanciation : on crée deux objets Voiture
ma_voiture = Voiture("rouge", "Ferrari")
ta_voiture = Voiture("bleue", "Renault")

ma_voiture.rouler()   # La Ferrari roule à 10 km/h
ta_voiture.rouler()   # La Renault roule à 10 km/h
\end{lstlisting}
  \begin{piege}
    \textbf{Différence clé :} \py{Voiture} est le \emph{plan} ;
    \py{ma\_voiture} est une voiture \emph{réelle et unique}.
    Chaque objet est indépendant, mais suit les règles de sa classe.
  \end{piege}
\end{frame}

% ================================================================= 3. constructeur
\section{Le constructeur}

\begin{frame}[fragile]{Le constructeur \py{\_\_init\_\_}}
  \begin{defi}[Constructeur]
    méthode spéciale appelée \textbf{automatiquement} à la création (instanciation)
    d'un objet. Elle initialise les attributs et prépare l'état initial de l'objet.
  \end{defi}
\begin{lstlisting}
class Voiture:
    def __init__(self, marque, modele, annee, couleur="blanc"):
        # Initialisation des attributs d'instance
        self.marque       = marque
        self.modele       = modele
        self.annee        = annee
        self.couleur      = couleur
        self.kilometrage  = 0        # Valeur par défaut
        self.est_demarree = False
\end{lstlisting}
\end{frame}

\begin{frame}[fragile]{Trois façons d'appeler le constructeur}
\begin{lstlisting}
# Création avec tous les paramètres
voiture1 = Voiture("Toyota", "Corolla", 2020, "rouge")

# Création avec paramètre optionnel par défaut
voiture2 = Voiture("Peugeot", "208", 2019)

# Création avec paramètres nommés
voiture3 = Voiture(marque="BMW", modele="X3", annee=2022, couleur="noir")
\end{lstlisting}
  \begin{methode}[À retenir]
    \py{self} désigne l'objet en cours de construction : il est passé
    automatiquement par Python, on ne le fournit jamais à l'appel.
  \end{methode}
\end{frame}

% ================================================================= 4. piliers
\section{Les trois piliers}

\begin{frame}[fragile,shrink=5]{Pilier 1 --- l'encapsulation}
  \begin{columns}[T,onlytextwidth]
    \begin{column}{0.5\textwidth}
      \textbf{Principe}
      \begin{itemize}
        \item on « cache » les détails internes ;
        \item on protège les données en contrôlant leur accès ;
        \item on évite les modifications accidentelles et on garantit la
              cohérence des données.
      \end{itemize}
      \smallskip
      \smallskip
      Attributs et méthodes sont soit \textbf{publics} (accessibles de
      l'extérieur), soit \textbf{privés} (accessibles seulement via des
      méthodes dédiées).
    \end{column}
    \begin{column}{0.48\textwidth}
      \begin{center}
      \begin{tikzpicture}[node distance=4mm]
        \node[classe] (cl) {Classe};
        \node[objet,below=of cl] (me) {Méthodes};
        \node[boite,below=of me,fill=qbg,draw=qbar] (va) {Variables};
        \node[boite,left=16mm of me,fill=cdcbg,draw=cdcbar,font=\small\bfseries] (en) {Encapsulation};
        \begin{scope}[on background layer]
          \node[draw=black!45,rounded corners=3pt,inner sep=4mm,fit=(cl)(me)(va)] (cap) {};
        \end{scope}
        \draw[fleche] (en.east) -- (cap.west);
      \end{tikzpicture}
      \end{center}
    \end{column}
  \end{columns}
  \vspace{-2mm}
  \begin{defi}[En Python, par convention]
    \py{solde} public --- \py{\_solde} protégé (usage interne) ---
    \py{\_\_solde} privé (accès via \py{deposer()}, \py{retirer()}).
  \end{defi}
\end{frame}

\begin{frame}[fragile]{Pilier 2 --- l'héritage}
  \begin{defi}[Héritage]
    une classe « enfant » hérite des attributs et méthodes d'une classe « parent ».
  \end{defi}
  \begin{columns}[T,onlytextwidth]
    \begin{column}{0.46\textwidth}
      \textbf{Avantages}
      \begin{itemize}
        \item réutilisation du code : rien à réécrire ;
        \item spécialisation : une \py{Voiture} \emph{est un} \py{Vehicule} ;
        \item l'enfant peut ajouter ses propres méthodes ou \textbf{surcharger}
              celles du parent.
      \end{itemize}
    \end{column}
    \begin{column}{0.52\textwidth}
\begin{lstlisting}
class Vehicule:                 # Parent
    def se_deplacer(self):
        print("Le véhicule se déplace.")

class Voiture(Vehicule):        # Enfant
    def se_deplacer(self):      # Surcharge
        print("La voiture roule.")
\end{lstlisting}
    \end{column}
  \end{columns}
\end{frame}

\begin{frame}[shrink=8]{Héritage : une hiérarchie de classes}
  \begin{center}
  \begin{tikzpicture}[
      every node/.style={boite,font=\footnotesize},
      level 1/.style={sibling distance=52mm},
      level 2/.style={sibling distance=18mm},
      level 3/.style={sibling distance=17mm},
      edge from parent/.style={draw=qbar,-{Stealth[length=2mm]},
        edge from parent path={(\tikzchildnode) -- (\tikzparentnode)}},
      level distance=13mm]
    \node[classe] {Être vivant}
      child { node[classe] {Animal}
        child { node {Insecte} }
        child { node {Mammifère}
          child { node {Chien} }
          child { node {Chauve-souris} } }
        child { node {Oiseau} } }
      child { node[classe] {Végétal}
        child { node {Algue} }
        child { node {Fleur} }
        child { node {Arbre}
          child { node {Pommier} }
          child { node {Cerisier} } } };
  \end{tikzpicture}
  \end{center}
  \vspace{-2mm}
  \begin{methode}[Lecture]
    la flèche se lit « \emph{est un} » : un \py{Pommier} est un \py{Arbre}, donc un
    \py{Végétal} --- et il hérite de tout ce qu'elles définissent.
  \end{methode}
\end{frame}

\begin{frame}[fragile,shrink=5]{Pilier 3 --- le polymorphisme}
  \begin{defi}[Polymorphisme]
    « plusieurs formes » : une même interface (une méthode de même nom) pour des
    objets de types différents. Chaque objet réagit à sa manière.
  \end{defi}
\begin{lstlisting}
class Voiture(Vehicule):
    def se_deplacer(self):
        print("La voiture roule.")

class Bateau(Vehicule):
    def se_deplacer(self):
        print("Le bateau navigue.")

for v in [Voiture(), Bateau()]:   # Même appel, comportements différents
    v.se_deplacer()
\end{lstlisting}
  \vspace{-2mm}
  \begin{methode}[Intérêt]
    le code appelant ignore le type exact : ajouter un \py{Avion} ne change
    rien à la boucle.
  \end{methode}
\end{frame}

% ================================================================= 5. compléments
\section{Compléments}

\begin{frame}[fragile]{Attributs de classe et attributs d'instance}
  \begin{columns}[T,onlytextwidth]
    \begin{column}{0.47\textwidth}
      \textbf{Attribut d'instance}
      \begin{itemize}
        \item propre à chaque objet ;
        \item défini avec \py{self.nom\_attribut} ;
        \item ex. le titulaire d'un compte.
      \end{itemize}
      \smallskip
      \textbf{Attribut de classe}
      \begin{itemize}
        \item partagé par tous les objets ;
        \item une seule copie en mémoire ;
        \item ex. le nom de la banque.
      \end{itemize}
    \end{column}
    \begin{column}{0.51\textwidth}
\begin{lstlisting}
class Compte:
    banque = "Banque Nationale"   # Attribut de classe

    def __init__(self, titulaire):
        self.titulaire = titulaire   # Attribut d'instance
\end{lstlisting}
      \smallskip
      \begin{piege}
        modifier \py{Compte.banque} change la valeur \emph{pour tous} les comptes.
      \end{piege}
    \end{column}
  \end{columns}
\end{frame}

\begin{frame}[fragile]{Accesseurs et mutateurs (\emph{getters} \& \emph{setters})}
  \begin{defi}[Rôle]
    donner un accès \textbf{contrôlé} à un attribut privé : l'accesseur
    (\emph{getter}) le lit, le mutateur (\emph{setter}) le modifie après vérification.
  \end{defi}
\begin{lstlisting}
class Compte:
    def __init__(self, titulaire, solde=0):
        self.titulaire = titulaire
        self.__solde   = solde          # Attribut privé

    def get_solde(self):                # Accesseur
        return self.__solde

    def deposer(self, montant):         # Mutateur contrôlé
        if montant <= 0:
            raise ValueError("Le montant doit être positif.")
        self.__solde += montant
\end{lstlisting}
\end{frame}

\begin{frame}[fragile,shrink=5]{La version pythonique : \py{@property}}
\begin{lstlisting}
class Compte:
    def __init__(self, solde=0):
        self.__solde = solde

    @property
    def solde(self):                # Lecture : compte.solde
        return self.__solde

    @solde.setter
    def solde(self, valeur):        # Écriture : compte.solde = 100
        if valeur < 0:
            raise ValueError("Le solde ne peut pas être négatif.")
        self.__solde = valeur
\end{lstlisting}
  \vspace{-2mm}
  \begin{methode}[Avantage]
    on écrit \py{compte.solde} comme un attribut ordinaire, mais le contrôle de
    cohérence reste assuré par la classe.
  \end{methode}
\end{frame}

% ================================================================= bilan
\begin{frame}{Bilan}
  \begin{center}
  \renewcommand{\arraystretch}{1.25}\small
  \begin{tabular}{@{}>{\bfseries}l l@{}}
    \toprule
    Notion & Définition en une ligne \\
    \midrule
    Classe          & le modèle : attributs + méthodes \\
    Objet           & une instance concrète de la classe \\
    Constructeur    & \py{\_\_init\_\_}, appelé à la création, initialise l'état \\
    Encapsulation   & cacher l'interne, contrôler l'accès aux données \\
    Héritage        & une classe enfant réutilise et spécialise un parent \\
    Polymorphisme   & un même appel, des comportements différents \\
    \bottomrule
  \end{tabular}
  \end{center}
  \vspace{-1mm}
  \begin{exam}
    \textbf{Suite :} TP POO Python --- écrire ses propres classes, encapsuler
    les attributs, construire une hiérarchie d'héritage.
  \end{exam}
\end{frame}

\end{document}
